% ============================================================================
% 07_draft_guards.tex  —  placeholder rendering and the final-build guard
%
% DRAFT MODE  (default):  every placeholder renders as a loud coloured box, and
%                         a tally is printed at the end of the log.
% FINAL MODE  (\finaltrue): the presence of ANY placeholder FAILS THE BUILD with
%                         an error naming every offender.
%
% Set final mode from main.tex or from the command line:
%     latexmk -pdf -usepretex="\finaltrue" main.tex
%
% The point of this file: a placeholder can only be removed by supplying the
% evidence behind it. Removing one without the evidence is the worst thing that
% can be done in this repository (CLAUDE.md §4).
% ============================================================================

\newif\iffinal\finalfalse

\definecolor{yhsaPending}{HTML}{C0181F}
\definecolor{yhsaNeeds}{HTML}{9B2FAE}
\definecolor{yhsaTodo}{HTML}{B8620A}

\newcounter{yhsaPendingCount}
\newcounter{yhsaNeedsCount}
\newcounter{yhsaTodoCount}
\newcounter{yhsaMissingNumCount}

\makeatletter

% --- the offender list, written to the log and to main.placeholders ----------
\newwrite\yhsa@phfile
\AtBeginDocument{\immediate\openout\yhsa@phfile=\jobname.placeholders\relax}
% #2 is DETOKENIZED before writing. Placeholder text legitimately contains
% macros such as \cref, and \immediate\write expands its argument -- expanding
% cleveref at write time, before the label exists, overflows the parameter stack.
\newcommand{\yhsa@record}[2]{%
  \immediate\write\yhsa@phfile{#1 | p.\thepage\space | \detokenize{#2}}%
  \iffinal
    \PackageError{yhsa-draft-guards}{%
      FINAL BUILD BLOCKED by an unresolved #1 on page \thepage: #2%
    }{%
      Every placeholder must be resolved before the final build.\MessageBreak
      A placeholder is resolved by SUPPLYING THE EVIDENCE, never by deleting the\MessageBreak
      marker. See CLAUDE.md section 4 and docs/DATA_REQUEST.md.\MessageBreak
      Run scripts/check_placeholders.py for the full list.%
    }%
  \fi
}

% --- \PENDING{key}{what is needed} : a missing NUMBER -----------------------
% INLINE and LINE-BREAKABLE. A missing number stands where the number will stand,
% inside a sentence. An \fbox there is an unbreakable box in mid-paragraph that
% overflows the margin and destroys the very page layout the draft build exists
% to show. Loud red is enough; a frame is not. The key is DETOKENIZED and set in
% \texttt because canonical keys contain underscores, which are subscript
% characters in text mode.
\newcommand{\yhsa@pending}[2]{%
  \stepcounter{yhsaPendingCount}%
  \yhsa@record{PENDING}{#1 -- #2}%
  \iffinal\else
    \allowbreak
    \textcolor{yhsaPending}{\textbf{[\,PENDING \texttt{\detokenize{#1}}: #2\,]}}%
  \fi
}

% --- a \num{} key with no definition ----------------------------------------
% Distinct from \PENDING: this means the key was USED but emit_numbers.py did
% not supply it, i.e. the canonical entry is still PENDING or the key is a typo.
\newcommand{\yhsa@missingnum}[1]{%
  \stepcounter{yhsaMissingNumCount}%
  \yhsa@record{MISSING-NUM}{\detokenize{#1}}%
  \iffinal\else
    \textcolor{yhsaPending}{\textbf{[?\,\texttt{\detokenize{#1}}]}}%
  \fi
}

% --- \NEEDSDATA{id}{dataset} : a missing FIGURE or TABLE --------------------
% Rendered as a box the size of the float it will become, so that the draft PDF
% shows the true shape and page budget of the finished report.
\newcommand{\yhsa@needsdata}[2]{%
  \stepcounter{yhsaNeedsCount}%
  \yhsa@record{NEEDSDATA}{#1 -- #2}%
  \iffinal\else
    \begin{center}
    \setlength{\fboxsep}{10pt}%
    \textcolor{yhsaNeeds}{%
      \fbox{\parbox[c][46mm][c]{\dimexpr0.88\linewidth-2\fboxsep-2\fboxrule\relax}{\centering\footnotesize\setstretch{1.0}%
        \textbf{\large NEEDS DATA}\\[4pt]
        \textbf{#1}\\[4pt]
        \textit{#2}\\[6pt]
        \tiny This box is a placeholder of approximately the size of the finished float.\\
        See docs/DATA\_REQUEST.md and figures/FIGURE\_REGISTER.md.}}%
    }
    \end{center}
  \fi
}

% --- \TODOPAL{question} : only Palaash can answer ---------------------------
% BLOCK-LEVEL. A question is a note about the surrounding text, not a word in it,
% so it is set as its own block. \par before and after keeps the framed box out
% of the enclosing paragraph, where it would overflow the right margin.
\newcommand{\yhsa@todopal}[1]{%
  \stepcounter{yhsaTodoCount}%
  \yhsa@record{TODOPAL}{#1}%
  \iffinal\else
    \par\medskip\noindent
    \textcolor{yhsaTodo}{%
      \fbox{\parbox{\dimexpr\linewidth-2\fboxsep-2\fboxrule\relax}{%
        \footnotesize\setstretch{1.0}\textbf{FOR PALAASH:} #1}}%
    }%
    \par\medskip
  \fi
}

% --- end-of-run tally --------------------------------------------------------
\AtEndDocument{%
  \immediate\closeout\yhsa@phfile
  \message{^^J}%
  \message{================ CHEM-151 PLACEHOLDER TALLY ================^^J}%
  \message{  PENDING      (missing number)        : \theyhsaPendingCount^^J}%
  \message{  MISSING-NUM  (undefined \string\num key)     : \theyhsaMissingNumCount^^J}%
  \message{  NEEDSDATA    (missing figure/table)  : \theyhsaNeedsCount^^J}%
  \message{  TODOPAL      (question for Palaash)  : \theyhsaTodoCount^^J}%
  \message{  full list: \jobname.placeholders^^J}%
  \message{============================================================^^J}%
}

\makeatother
